\documentclass[openany,a4paper,12pt]{book}

% --- Lingua e Font ---
\usepackage[utf8]{inputenc}
\usepackage[T1]{fontenc} 
\usepackage[italian]{babel}
\usepackage{lmodern} % Font più nitido in PDF rispetto al default
\usepackage{microtype} % Migliora spaziature e sillabazione

% --- Colori (Definiti prima per usarli ovunque) ---
\usepackage[x11names, dvipsnames]{xcolor}
\definecolor{PoliBlue}{RGB}{0, 45, 114} % Blu accademico elegante
\definecolor{OpticRed}{RGB}{200, 10, 30} % Rosso "laser" per i link
\definecolor{ForestGreen}{RGB}{34, 139, 34}
\definecolor{CodeBack}{RGB}{245, 245, 245}
\definecolor{BoxGray}{RGB}{248, 249, 250}

% --- Matematica e Fisica ---
\usepackage{amsmath, amssymb, amsfonts, amsthm}
\usepackage{mathtools, cancel, bm}

% Non usare physics insieme a siunitx: i due pacchetti definiscono entrambi \qty e generano conflitti.
% Per le unità di misura usa solo siunitx.
\usepackage{siunitx}
\sisetup{
    output-decimal-marker={,},
    separate-uncertainty=true, % Per scrivere (10,5 ± 0,1) m
    per-mode=symbol, % Scrive m/s invece di m s^{-1}
    inter-unit-product = \ensuremath{{}\cdot{}} % Punto tra le unità
}

% --- Formattazione Testo e Layout ---
\usepackage{ragged2e}
\usepackage{quoting}
\usepackage{booktabs} % Per tabelle professionali
\usepackage{enumitem}
\usepackage{array, tabularx}
\usepackage[a4paper, margin=2.5cm, headheight=15pt]{geometry} 
\setlength{\parindent}{0pt}
\setlength{\parskip}{0.75em}

% --- Titoli di Sezione (Stile Libro Universitario Moderno) ---
\usepackage{titlesec}
\titleformat{\chapter}[display]
  {\normalfont\bfseries\color{PoliBlue}}{\filleft\Huge\thechapter}{2ex}
  {\titlerule\vspace{2ex}\Huge\filleft}
  [\vspace{2ex}\titlerule]
\titleformat{\section}
  {\normalfont\Large\bfseries\color{PoliBlue}}{\thesection}{1em}{}
\titleformat{\subsection}
  {\normalfont\large\bfseries\color{PoliBlue}}{\thesubsection}{1em}{}

% --- Grafica, TikZ e Plot ---
\usepackage{graphicx}
\usepackage{float} 
\usepackage{circuitikz}
\usepackage{pgfplots}
\pgfplotsset{compat=1.18}
\usepackage{subcaption} 
\usepackage{tikz}
\usepackage{tikz-3dplot}
\usetikzlibrary{positioning, arrows.meta, angles, quotes, trees, patterns, calc}

\pgfplotsset{
    standard/.style={
        axis lines=middle,
        xlabel={$t$}, ylabel={$y$},
        grid=both,
        grid style={line width=.1pt, draw=gray!20},
        major grid style={line width=.2pt, draw=gray!40},
        samples=100,
        enlargelimits=true,
        label style={font=\small},
        tick label style={font=\tiny}
    }
}

% --- Box ed Evidenziazioni (Tcolorbox) ---
\usepackage[most]{tcolorbox}

% Stile per le Note (si spezzano se a fine pagina)
\newtcolorbox{nota}[1][]{
  enhanced, breakable,
  left=1em, right=1em, top=1ex, bottom=1ex,
  colback=BoxGray, colframe=gray!50,
  borderline west={3pt}{0pt}{gray},
  fontupper=\itshape,
  title={\textbf{Nota}}, coltitle=black, attach title to upper={\quad},
  #1
}

% Stile per le Definizioni
\newtcolorbox{definizione}[1][]{
  enhanced, breakable,
  left=1em, right=1em, top=1ex, bottom=1ex,
  colback=PoliBlue!5, colframe=PoliBlue!80,
  borderline west={3pt}{0pt}{PoliBlue},
  title={\textbf{Definizione}}, coltitle=PoliBlue!80!black, attach title to upper={\quad},
  #1
}

% Stile per i Teoremi/Formule importanti
\newtcbtheorem[no counter]{teorema}{Teorema}%
{enhanced, breakable, colback=OpticRed!5, colframe=OpticRed, fonttitle=\bfseries}{th}

% --- Configurazione LISTINGS (Python per Analisi Dati) ---
\usepackage{listings}
\lstset{
    language=Python,
    backgroundcolor=\color{CodeBack},
    basicstyle=\ttfamily\small,
    keywordstyle=\color{PoliBlue}\bfseries,
    stringstyle=\color{OpticRed},
    commentstyle=\color{ForestGreen}\itshape,
    numberstyle=\tiny\color{gray},
    numbers=left,
    numbersep=10pt,
    frame=leftline,
    framerule=2pt,
    rulecolor=\color{PoliBlue!50},
    showspaces=false,
    showstringspaces=false,
    breaklines=true,
    breakatwhitespace=false,
    columns=fullflexible, % Permette un copia-incolla perfetto dal PDF
    keepspaces=true,
    literate={à}{{\`a}}1 {è}{{\`e}}1 {ì}{{\`i}}1 {ò}{{\`o}}1 {ù}{{\`u}}1 {±}{{$\pm$}}1
}

% --- Intestazioni e Piè di Pagina ---
\usepackage{fancyhdr}
\pagestyle{fancy}
\fancyhf{}
\renewcommand{\headrulewidth}{0.4pt}
\renewcommand{\chaptermark}[1]{\markboth{\thechapter.\ #1}{}}
\renewcommand{\sectionmark}[1]{\markright{\thesection.\ #1}}
\fancyhead[LE]{\small\textbf{\nouppercase{\leftmark}}} % Capitolo a sx (pagine pari)
\fancyhead[RO]{\small\textbf{\nouppercase{\rightmark}}} % Sezione a dx (pagine dispari)
\fancyfoot[C]{\small\textbf{\thepage}}

\fancypagestyle{plain}{%
    \fancyhf{}%
    \renewcommand{\headrulewidth}{0pt}%
    \fancyfoot[C]{\small\textbf{\thepage}}%
}

% --- Link e Indici (DEVE ESSERE L'ULTIMO PACCHETTO) ---
\usepackage{hyperref}
\usepackage{bookmark}
\hypersetup{
    colorlinks=true,
    linkcolor=PoliBlue, 
    urlcolor=OpticRed,
    citecolor=ForestGreen,
    pdftitle={Appunti di Ottica},
    pdfauthor={Mirolo, Brusini, Tedeschini, Stroili}
}

% ==================================================
% INIZIO DOCUMENTO
% ==================================================
\begin{document}

% --- FRONTESPIZIO CUSTOM ---
\begin{titlepage}
    \centering
    
    % Spazio iniziale dal margine superiore
    \vspace*{3cm}
    

    \vspace{0.8cm}
    
    % Titolo Principale
    {\Huge \bfseries \textcolor{PoliBlue}{Appunti del Laboratorio di Ottica} \par}
    
    \vspace{0.6cm}
    
    % Spazio centrale
    \vspace{3.5cm}
    
    \vspace{1.5cm}
    
    % Tabella autori con spazio aggiuntivo tra le colonne e le righe
    {
    \large
    \renewcommand{\arraystretch}{1.5} % Aumenta lo spazio verticale tra i nomi
    \begin{tabular}{c @{\hspace{2.5cm}} c} % @{\hspace{...}} allontana le due colonne
         \textbf{Mirolo Manuele} & \textbf{Alessio Brusini} \\
         \textbf{Michele Tedeschini} & \textbf{Emanuele Stroili}
    \end{tabular}
    \par
    }
    
    % \vfill spinge tutto quello che c'è sotto verso il fondo della pagina
    \vfill
    
    % Anno Accademico a fondo pagina
    {\large Anno Accademico 2026/2027 \par}
    
    % Piccolo spazio finale per non attaccarlo troppo al bordo del foglio
    \vspace*{1.5cm}
\end{titlepage}

\justifying

\frontmatter % Numerazione romana per indici
\tableofcontents

\mainmatter % Numerazione araba per i capitoli
\pagestyle{fancy}

\chapter{Introduzione all'Ottica}
Benvenuti nel laboratorio di ottica. In questo capitolo mostriamo le funzionalità del template.

\section{Esempi di Formattazione}

Un esperimento tipico richiede la misura della lunghezza d'onda di un laser He-Ne. Secondo i nostri dati, il valore misurato è $\lambda = \qty{632.8 \pm 0.1}{\nano\metre}$. 

Per scrivere un'espressione con parentesi in modo standard, si può usare:
\begin{equation}
    f(x) = \sin\left(\frac{2\pi}{\lambda} x\right)
\end{equation}

\begin{definizione}
L'ottica geometrica è l'approssimazione in cui la lunghezza d'onda della luce è considerata molto minore delle dimensioni degli oggetti con cui interagisce ($\lambda \to 0$).
\end{definizione}

\begin{nota}
Ricordarsi di allineare sempre gli specchi dell'interferometro di Michelson prima di iniziare la presa dati!
\end{nota}

\begin{teorema}{Legge di Snell}{snell}
Dati due mezzi con indici di rifrazione $n_1$ e $n_2$, l'angolo di incidenza $\theta_1$ e l'angolo di rifrazione $\theta_2$ sono legati da:
\begin{equation}
    n_1 \sin\theta_1 = n_2 \sin\theta_2
\end{equation}
\end{teorema}

\subsection{Analisi Dati}
Ecco un esempio di script per il fit dei dati d'interferenza:

\begin{lstlisting}
import numpy as np
import matplotlib.pyplot as plt

# Dati dell'esperimento di diffrazione
posizioni = np.array([0.1, 0.2, 0.3]) # ± 0.01 mm
intensita = np.array([100, 50, 20])

plt.plot(posizioni, intensita, 'ro-')
plt.title("Figura di diffrazione")
plt.xlabel("Posizione [mm]")
plt.ylabel("Intensità [u.a.]")
plt.grid(True)
plt.show()
\end{lstlisting}

\end{document}